\chapter{The Competitive Industry}\label{ch:b1:competitive-industry}

In March 2014 a licence to run a taxi in New York City, a medallion, sold for
an average of 1.16 million dollars (for the corporate, wheelchair-accessible
kind); in November 2019 the average medallion sold for
164,518 dollars. % ledger: b1:nyc-medallions
The number of medallions was capped at about 13,600. The cap kept the earnings of each cab, at fares set by the city, above
what \omterm{def:b1:competitive-industry:free-entry}{free entry} would have allowed, and the
price of a medallion was the value of those extra earnings. When ride-hailing
apps brought in thousands of cars without medallions, the extra earnings,
and the price, collapsed. This chapter builds an industry from the firms of
\cref{ch:b1:profit-maximisation}, lets firms enter and leave, and follows
the payments that entry cannot remove --- rents --- into the price of the
assets that earn them.

\begin{openingchart}{Real house prices, five countries, 1870--2020 (Canada from
1921, the United Kingdom from 1899, the United States from 1890), nominal
house price index divided by the consumer price index, 1990 = 100,
logarithmic axis.}{Jord\`a--Schularick--Taylor Macrohistory Database: house
prices (nominal) and consumer prices}
\begin{tikzpicture}
\begin{semilogyaxis}[omtime, log ticks with fixed point, ytick={10,20,50,100,200}, xmin=1870, xmax=2020,
  ymin=6, ymax=320, xlabel={Year}, ylabel={Real house price (1990 = 100)},
  width=0.72\linewidth, height=0.46\linewidth, legend pos=south east, unbounded coords=jump]
  \addplot+[no marks, thick] table[x=year, y=aus, col sep=comma] {figdata/bachelor-1/12-competitive-industry/out/real_house_prices.csv};
  \addplot+[no marks, thick] table[x=year, y=can, col sep=comma] {figdata/bachelor-1/12-competitive-industry/out/real_house_prices.csv};
  \addplot+[no marks, thick] table[x=year, y=fra, col sep=comma] {figdata/bachelor-1/12-competitive-industry/out/real_house_prices.csv};
  \addplot+[no marks, thick] table[x=year, y=gbr, col sep=comma] {figdata/bachelor-1/12-competitive-industry/out/real_house_prices.csv};
  \addplot+[no marks, thick] table[x=year, y=usa, col sep=comma] {figdata/bachelor-1/12-competitive-industry/out/real_house_prices.csv};
  \legend{Australia, Canada, France, United Kingdom, United States}
\end{semilogyaxis}
\end{tikzpicture}
\end{openingchart}

For eighty years real house prices went nowhere in particular; after 1950
they rose several times over. From 1950 to 2020 they were multiplied by 5.1
in the United Kingdom, 5.3 in Australia and 22.5 in France (from a very low
point after the wars), and by 1.9 in the United States. % figdata: house_points
Builders are a competitive industry with \omterm{def:b1:competitive-industry:free-entry}{free entry}, so a house cannot cost
much more than it costs to build for long. What rose was the price of the
land under it: a study of fourteen countries attributes about $80\%$ of the
rise in house prices since the Second World War to land (Knoll, Schularick
and Steger 2017). % ledger: b1:knoll2017
Land in the places people want to live is in fixed supply, and fixed
supply earns a rent.

\begin{recall}
In a competitive market firms enter when economic profit is positive and
leave when it is negative; the present value of $R$ a year for ever at the
interest rate $i$ is $R/i$. Both were met at school. A firm's supply is its
marginal cost above the \omterm{def:b1:profit-maximisation:shut-down}{shut-down price}
(\cref{ch:b1:profit-maximisation}).
\end{recall}

\section{The short run: a given number of firms}\label{sec:b1:competitive-industry:short-run}

\begin{definition}[Industry supply]\label{def:b1:competitive-industry:industry-supply}
\index{industry supply}\index{supply curve!of an industry}
With $n$ firms whose supply functions are $q_i(p)$, the short-run
\emph{industry supply} is their sum at each price,
$Q(p) = \sum_{i=1}^n q_i(p)$: the horizontal sum of their supply curves.
\end{definition}

Brenland's fishing boats (invented) each have the \omterm{def:b1:cost-minimisation:cost-function}{cost function}
$C(q) = 50 + 2q^2$, in thousands of dalens a year for $q$ tonnes of fish. A
boat's marginal cost is $4q$, so it supplies $q = p/4$, and 400 boats supply
$Q = 100p$. Against the demand $Q = 3000 - 50p$, the price is 20 thousand
dalens a tonne, each boat lands 5 tonnes and the industry 2,000. The
short-run \omterm{def:b1:competitive-industry:industry-supply}{industry supply} is flatter than one firm's supply curve: a rise
in price of one thousand dalens adds a quarter of a tonne for each boat, 100
tonnes for the industry.

\section{The long run: free entry}\label{sec:b1:competitive-industry:long-run}

\begin{definition}[Free entry]\label{def:b1:competitive-industry:free-entry}
\index{free entry}\index{entry and exit}
An industry has \emph{free entry} if any number of firms with access to the
same technology and input prices can enter it, and incumbents can leave,
without paying a cost that those already in it did not pay.
\end{definition}

With \omterm{def:b1:competitive-industry:free-entry}{free entry} and identical firms, the long-run price cannot stay above
the minimum of average cost, or firms would enter and make a profit; it
cannot stay below it, or firms would leave. So $p = \min AC$, each firm
produces at its \omterm{def:b1:cost-minimisation:mes}{minimum efficient scale}, and the number of firms is
whatever it takes to meet demand at that price. For the boats, average
cost $50/q + 2q$ is lowest at $q = 5$, where it is 20: with 2,000 tonnes
demanded at that price, the industry holds 400 boats.

\begin{definition}[Long-run industry supply]\label{def:b1:competitive-industry:lr-supply}
\index{long-run industry supply}
The \emph{long-run \omterm{def:b1:competitive-industry:industry-supply}{industry supply}} gives the quantity the industry
offers at each price once the number of firms has adjusted to make profit
zero.
\end{definition}

\begin{definition}[Constant-cost and increasing-cost industries]\label{def:b1:competitive-industry:cost-industry}
\index{constant-cost industry}\index{increasing-cost industry}
An industry is \emph{constant-cost} if its expansion leaves the input prices
its firms pay unchanged, so its long-run supply is flat at the minimum of
average cost; it is \emph{increasing-cost} if its expansion raises the price
of an input in limited supply, so its long-run supply rises.
\end{definition}

\begin{omfigure}
\begin{tikzpicture}
\begin{groupplot}[group style={group size=2 by 1, horizontal sep=16mm}, omrecord, grid=none,
  width=0.44\linewidth, height=0.42\linewidth, samples=80, ymin=0, ymax=40]
\nextgroupplot[xlabel={Boat's catch $q$}, ylabel={Thousand dalens a tonne}, xmin=0, xmax=10,
  title={\small One boat}, width=0.44\linewidth]
  \fill[omDef!14] (axis cs:0,20.24) rectangle (axis cs:5.83,23.33);
  \addplot[shifted, domain=1.5:10] {50/x + 2*x};
  \addplot[supply, domain=0:10] {4*x};
  \addplot[guide] coordinates {(0,20) (10,20)};
  \addplot[guide] coordinates {(0,23.33) (10,23.33)};
  \addplot[only marks, mark=*, black] coordinates {(5,20) (5.83,23.33)};
  \node[font=\footnotesize, anchor=north west] at (axis cs:6.5,38) {$MC$};
  \node[font=\footnotesize, anchor=south] at (axis cs:2,30) {$AC$};
\nextgroupplot[xlabel={Industry catch $Q$}, xmin=0, xmax=3500, title={\small The industry},
  width=0.44\linewidth, scaled x ticks=false, xticklabel style={/pgf/number format/.cd, fixed, 1000 sep={\,}}]
  \addplot[demand, domain=0:3000] {60 - x/50};
  \addplot[demand, dashed, domain=0:3500] {70 - x/50};
  \addplot[supply, domain=0:3500] {x/100};
  \addplot[supply, dashed, domain=0:3500] {x/125};
  \addplot[shifted] coordinates {(0,20) (3500,20)};
  \addplot[only marks, mark=*, black] coordinates {(2000,20) (2333,23.33) (2500,20)};
\end{groupplot}
\end{tikzpicture}

{\small Demand for fish rises from $Q = 3000 - 50p$ to $3500 - 50p$ (blue, dashed).
With 400 boats (red) the price rises to 23.33 and each boat lands 5.83
tonnes at an average cost of 20.24: the shaded profit. Boats enter until
there are 500 (red, dashed), and the price returns to 20, the long-run
supply (black, dashed) of a \omterm{def:b1:competitive-industry:cost-industry}{constant-cost industry}.}
\end{omfigure}

In a \omterm{def:b1:competitive-industry:cost-industry}{constant-cost industry}, a rise in demand raises the price only for as
long as entry takes; the long-run effect falls on the number of firms.
Most industries draw on some input whose supply is limited --- fishing
grounds, land near cities, skilled people --- and then the long-run supply
rises. If the industry's long-run supply is $p = 10 + Q/100$ because each
extra tonne must come from poorer or more distant grounds, and demand is
$Q = 2000 - 50p$, the price is 20 and the quantity 1,000. Every tonne but the
last costs less than 20 to land; the difference goes to whoever owns the
better grounds.

\section{Rents and quasi-rents}\label{sec:b1:competitive-industry:rents}

\begin{definition}[Economic rent]\label{def:b1:competitive-industry:rent}
\index{economic rent}\index{rent (economic)}
An \emph{economic rent} is a payment to a factor of production above the
least payment that would keep it in its current use. A factor in fixed
supply earns a pure rent: all of its payment is rent, since it would be
supplied at any price.
\end{definition}

In the increasing-cost example the owners of the fishing grounds earn the
triangle between the price and the long-run supply curve,
$\tfrac12 \times 1000 \times (20 - 10) = 5000$ thousand dalens a year. Competition does not
remove it: entrants would have to use worse grounds, at a cost of 20. Any
profit a firm seems to make from a better ground is a rent to the ground,
and if the firm rents it, competition among firms passes the rent to the
owner.

\begin{omfigure}
\begin{tikzpicture}
\begin{axis}[omrecord, grid=none, xlabel={Quantity $Q$}, ylabel={Price}, xmin=0, xmax=2100, ymin=0, ymax=45,
  width=0.52\linewidth, height=0.4\linewidth, samples=60, scaled x ticks=false,
  xticklabel style={/pgf/number format/.cd, fixed, 1000 sep={\,}}]
  \addplot[cs area, forget plot] coordinates {(0,10) (1000,20) (0,20)} -- cycle;
  \addplot[demand, domain=0:2000] {40 - x/50};
  \addplot[supply, domain=0:2000] {10 + x/100};
  \addplot[guide] coordinates {(0,20) (1000,20) (1000,0)};
  \node[font=\footnotesize] at (axis cs:300,17.3) {rent};
\end{axis}
\end{tikzpicture}

{\small An \omterm{def:b1:competitive-industry:cost-industry}{increasing-cost industry}: long-run supply $p = 10 + Q/100$ (red),
demand $Q = 2000 - 50p$ (blue). The shaded triangle, 5,000 a year, is the rent
earned by the owners of the scarce input.}
\end{omfigure}

\begin{definition}[Quasi-rent]\label{def:b1:competitive-industry:quasi-rent}
\index{quasi-rent}
A \emph{quasi-rent} is revenue minus variable cost: the payment to the
fixed factors of a firm in the short run. It is a rent while the factors
are fixed and becomes a cost when they must be replaced.
\end{definition}

A boat that lands 5.83 tonnes at the price 23.33 earns a \omterm{def:b1:competitive-industry:quasi-rent}{quasi-rent} of
$23.33 \times 5.83 - 2 \times 5.83^2 \approx 68.04$ thousand dalens (68.06 unrounded), against its fixed
cost of 50. A firm stays in business in the short run if its \omterm{def:b1:competitive-industry:quasi-rent}{quasi-rent} is
positive, and stays in the long run only if it covers the fixed cost.

\section{Capitalisation}\label{sec:b1:competitive-industry:capitalisation}

\begin{definition}[Capitalisation]\label{def:b1:competitive-industry:capitalisation}
\index{capitalisation}
\emph{Capitalisation} is the inclusion of an expected stream of rents in
the price of the asset that earns them: with free trade in the asset, its
price is the present value of the rents, $R/i$ for a permanent rent $R$ at
the interest rate $i$.
\end{definition}

In the capped fishery of the figure below, a licence lets its owner earn
62,500 dalens a year above the normal return. It sells, at $4\%$, for $62500/0.04 = 1.5625$ million. The person who
received it free when the cap was set gains that sum. Whoever buys it later
pays the price and earns only $4\%$ on it: the normal return. The rent has
become a cost, and the buyer's earnings look ordinary, although the cap
still keeps prices up for the industry's customers. When the rule changes,
the price of the licence falls at once, by the value of the lost rents,
and the loss falls on whoever owns it then.

\begin{omfigure}
\begin{tikzpicture}
\begin{axis}[omrecord, grid=none, xlabel={Catch $Q$ (tonnes)}, ylabel={Thousand dalens a tonne},
  xmin=0, xmax=3100, ymin=0, ymax=62, width=0.6\linewidth, height=0.44\linewidth, samples=60,
  scaled x ticks=false, xticklabel style={/pgf/number format/.cd, fixed, 1000 sep={\,}}]
  \fill[omDef!20] (axis cs:0,21.67) rectangle (axis cs:1500,30);
  \fill[omThm!25] (axis cs:0,20) rectangle (axis cs:1500,21.67);
  \fill[omProp!25] (axis cs:1500,30) -- (axis cs:2000,20) -- (axis cs:1500,20) -- cycle;
  \addplot[demand, domain=0:3000] {60 - x/50};
  \addplot[supply, domain=0:3000] {x/50};
  \addplot[shifted] coordinates {(0,20) (3000,20)};
  \addplot[only marks, mark=*, black] coordinates {(1500,30) (2000,20)};
  \node[font=\footnotesize] at (axis cs:750,25.8) {licence rents};
  \node[font=\footnotesize, anchor=west] at (axis cs:2450,50) {200 boats};
  \node[font=\footnotesize, anchor=south] at (axis cs:2600,20) {free entry};
\end{axis}
\end{tikzpicture}

{\small A cap of 200 boats: their supply $Q = 50p$ (red) meets demand at a price of 30
and a catch of 1,500 tonnes. The top rectangle is the licence rents,
$200 \times 62.5 = 12500$ thousand dalens a year; the thin rectangle is the extra cost
of boats fishing above their \omterm{def:b1:cost-minimisation:mes}{minimum efficient scale}; the triangle is the
catch not landed. The last two, 5,000 a year, are what the cap costs.}
\end{omfigure}

\begin{record}[Taxi medallions]
New York City limits the number of taxi medallions, about 13,600 in 2019.
The average price of a corporate wheelchair-accessible medallion was 1.16
million dollars in March 2014; the average price of all medallions sold in
November 2019 was 164,518 dollars. The City Council attributes the fall to
large-scale disruption of the industry by app-based companies, which
could put cars on the street without medallions. % ledger: b1:nyc-medallions
\omterm{def:b1:competitive-industry:capitalisation}{Capitalisation} explains the high price, the value of the rents that a
capped licence was expected to earn, and why the price fell as soon as
those rents were expected to shrink. The loss fell on whoever owned a
medallion when the fall came, including those who had just paid the peak
price.
\end{record}

\begin{model}{Free entry}
\assumes Many potential firms with the same U-shaped average cost, free to
enter and leave, at given input prices.
\predicts Zero long-run profit, a price at the minimum of average cost, and
entry after a rise in demand that brings the price back down; a flat
long-run supply unless some input is scarce.
\tested The post-war rise in house prices was mostly a rise in the price of
land, the input in fixed supply, not in the cost of building: about $80\%$
of it in Knoll, Schularick and Steger (2017). % ledger: b1:knoll2017
\breaks Licences and patents block entry; when firms differ in cost, the
low-cost firms keep inframarginal rents, which then belong to whatever
makes them cheaper.
\end{model}

\begin{model}{Capitalisation}
\assumes A known, permanent stream of rents, a discount rate, and an asset
that can be bought and sold.
\predicts The asset's price equals the present value of the rents; a rule
change is priced in at once; later buyers earn the normal return.
\tested Consistent with it: taxi medallion prices were high while the cap
held and fell steeply once app-based cars entered. % ledger: b1:nyc-medallions
A test would compare the fall in an asset's price with an independent
estimate of the present value of the rents lost, measured when the rule
changed.
\breaks When the stream is uncertain, when buyers are wrong about it, or
when the asset cannot be sold (a licence tied to its owner), the price need
not follow the present value.
\end{model}

\begin{measure}[A house-price index]
A house-price index compares the prices of houses sold in different
years. Since no two houses are alike, the index must correct for what is
sold: \emph{repeat-sales} indexes use the change in price of the same house
sold twice; \emph{hedonic} indexes regress prices on the characteristics of
the houses (size, rooms, location) and track the price of a fixed bundle of
characteristics. Long series join several such indexes and older sources,
so their early decades are less precise.
\end{measure}

\section*{Exercises}\label{sec:b1:competitive-industry:exercises}

\begin{exercise}[$\star$]\label{exo:b1:competitive-industry:sum}
Fifty identical firms each supply $q = p/2$. Write the \omterm{def:b1:competitive-industry:industry-supply}{industry supply}.
\end{exercise}

\begin{exercise}[$\star$]\label{exo:b1:competitive-industry:entry}
Identical firms have $C(q) = 100 + q^2$, entry is free and demand is
$Q = 1000 - 10p$. Find the long-run price, the output of each firm and the
number of firms.
\end{exercise}

\begin{exercise}[$\star$]\label{exo:b1:competitive-industry:read}
From the opening chart's table, the real house price index was 35.8 in
1950 and 181.9 in 2020 in the United Kingdom, and 84.8 and 158.5 in the
United States. % figdata: house_points
Compute the ratios.
\end{exercise}

\begin{exercise}[$\star$]\label{exo:b1:competitive-industry:demand-rise}
Demand rises by $20\%$ at every price in a \omterm{def:b1:competitive-industry:cost-industry}{constant-cost industry}. Predict
the price and the number of firms in the short run and in the long run.
\end{exercise}

\begin{exercise}[$\star\star$]\label{exo:b1:competitive-industry:increasing}
An industry's long-run supply is $p = 10 + Q/100$ and demand is $Q = 2000 - 50p$.
Find the price, the quantity and the rent to the scarce input. Demand rises
to $Q = 2600 - 50p$: what happens to the rent?
\end{exercise}

\begin{exercise}[$\star\star$]\label{exo:b1:competitive-industry:quasi}
A firm with $C(q) = 50 + 2q^2$ faces the price 24. Compute its output, its
\omterm{def:b1:competitive-industry:quasi-rent}{quasi-rent} and its profit.
\end{exercise}

\begin{exercise}[$\star\star$]\label{exo:b1:competitive-industry:licence}
A licence lets its owner earn 40,000 dollars a year above the normal return,
for ever. What is it worth at an interest rate of $5\%$?
\end{exercise}

\begin{exercise}[$\star\star$]\label{exo:b1:competitive-industry:more-licences}
In the weekend problem's fishery the authority raises the number of
licences from 200 to 240. Compute the new price, the rent per licence and
the licence price at $4\%$. By what percentage does the licence price fall?
\end{exercise}

\begin{exercise}[$\star\star$]\label{exo:b1:competitive-industry:growth}
Using the ratios of \cref{exo:b1:competitive-industry:read}, compute the
yearly growth of real house prices from 1950 to 2020 in both countries.
\end{exercise}

\begin{exercise}[$\star\star\star$]\label{exo:b1:competitive-industry:subsidy}
A subsidy of $s$ per unit is given to the producers of a competitive
industry. Show that in the long run buyers gain all of it if the industry
is constant-cost, and that in an \omterm{def:b1:competitive-industry:cost-industry}{increasing-cost industry} part of it goes
to the owners of the scarce input.
\end{exercise}

\begin{exercise}[$\star\star\star$]\label{exo:b1:competitive-industry:expected}
A city announces that its taxi cap will be lifted in ten years. Who loses:
the owners at the announcement, or those who own medallions in ten years?
Compare with a cap lifted without warning.
\end{exercise}

\begin{exercise}[$\star\star\star$]\label{exo:b1:competitive-industry:refute}
``House prices rose because land is scarce.'' What observation would refute
this? Name one comparison between places and one over time.
\end{exercise}

\section{Problem: Fishing Licences in Brenland}\label{sec:b1:competitive-industry:problem}

\begin{problem}[{Weekend problem --- the capitalised value of a licence}]\label{pb:b1:competitive-industry:fishing}
Each Brenland fishing boat (invented) has the \omterm{def:b1:cost-minimisation:cost-function}{cost function}
$C(q) = 50 + 2q^2$, in thousands of dalens a year, where $q$ is its catch in
tonnes; the fixed cost includes a normal return on the boat. Demand for
fish is $Q = 3000 - 50p$, with $p$ in thousands of dalens a tonne.

\medskip\noindent\textbf{Part I --- Free entry.}
\begin{enumerate}
  \item Compute marginal and average cost, and the minimum of average cost.
  \item Write one boat's supply.
  \item What is the long-run price with \omterm{def:b1:competitive-industry:free-entry}{free entry}?
  \item Compute the catch.
  \item How many boats fish?
  \item Demand rises to $Q = 3500 - 50p$. Find the long-run price and the number of boats.
\end{enumerate}

\medskip\noindent\textbf{Part II --- A cap.} Before demand rose, the authority
caps the fleet at 200 boats, with one licence each.
\begin{enumerate}[resume]
  \item Write the \omterm{def:b1:competitive-industry:industry-supply}{industry supply} of 200 boats.
  \item Find the price and the catch.
  \item Compute each boat's catch and average cost.
  \item Compute each boat's profit.
  \item Compute each boat's \omterm{def:b1:competitive-industry:quasi-rent}{quasi-rent}.
  \item Why does entry not remove the profit?
  \item What is the rent earned by a licence?
\end{enumerate}

\medskip\noindent\textbf{Part III --- Capitalisation.} Licences can be sold;
the interest rate is $4\%$.
\begin{enumerate}[resume]
  \item What is a licence worth?
  \item A fisher buys a licence. What return does she earn on what she paid?
  \item Who gained from the cap?
  \item What is a licence worth if the interest rate falls to $2.5\%$?
  \item The authority announces that the cap will be lifted in ten years. What is a licence worth now?
  \item How much did the announcement cost the owner of one licence?
\end{enumerate}

\medskip\noindent\textbf{Part IV --- The cap goes.} The cap is lifted without warning.
\begin{enumerate}[resume]
  \item Find the new price and the number of boats.
  \item By how much does the price of a licence fall?
  \item Compute \omterm{def:b1:consumer-welfare:consumer-surplus}{consumer surplus} with the cap and without it.
  \item Show that buyers gain more than licence owners lose, and explain the difference.
  \item Who bears the loss?
  \item State the result: the capitalised value of one licence, and its fall when the cap goes.
\end{enumerate}
\end{problem}
